\chapter{Sprint 5: Frontend Integration and Live Monitoring}
\chaptertoc

\clearpage

\section{Introduction}
Sprint 5 connected the dashboard to the live backend services and runtime events. At this stage, the frontend stopped being only a visual shell and became a monitoring interface that could react to feed changes and live metrics.

\section{Sprint Backlog}
The Sprint 5 backlog focused on the pieces needed to make the dashboard react to live data:
\begin{itemize}
    \item connect the dashboard to backend responses in a reliable way,
    \item allow the user to manage feeds from the interface,
    \item receive live updates without refreshing the page,
    \item choose the most suitable video transport depending on availability,
    \item keep the dashboard usable when the live connection is interrupted,
    \item validate the main setup and onboarding flows through integration tests.
\end{itemize}

\section{Design}

\subsection{Data Flow Diagram}
\begin{figure}[H]
    \centering
    \fbox{\parbox[c][6cm][c]{0.85\textwidth}{\centering Placeholder for Sprint 5 frontend data flow diagram}}
    \caption{Data flow between backend services and frontend monitoring components.}
    \label{fig:sprint5_data_flow}
\end{figure}

In practice, the backend handles configuration and diagnostic requests, live events carry the real-time updates, and the dashboard groups those signals into KPI cards, activity logs, and attention panels.

\subsection{Sequence Diagrams}
\begin{figure}[H]
    \centering
    \fbox{\parbox[c][6cm][c]{0.85\textwidth}{\centering Placeholder for Sprint 5 sequence diagram}}
    \caption{Sprint 5 sequence diagram for dashboard updates and alert surfacing.}
    \label{fig:sprint5_sequence}
\end{figure}

Three main flows were implemented:
\begin{enumerate}
    \item \textbf{Feed lifecycle:} a dashboard action updates the feed state and the new status is sent back to the interface.
    \item \textbf{Live metrics:} metric updates are received from the backend and reflected immediately on the dashboard.
    \item \textbf{Alert surfacing:} alerts reach the interface and appear in the activity and attention areas.
\end{enumerate}

\section{Implementation}
\subsection{API Integration}
The frontend relies on a dedicated API layer that keeps the communication with the backend structured. It defines the main data types used by the dashboard and centralizes request handling so errors can be managed consistently.

\subsection{Authentication and Role-Based Navigation}
Sprint 5 also connected the login and sign-up flows to the dashboard. The auth provider clears cached data when the session changes, protected routes redirect users according to their role, and administrators are taken to the manager administration page while managers are routed to the monitoring dashboard. This keeps the frontend aligned with the backend access rules and avoids exposing the wrong interface to the wrong user.

\subsection{Dashboard State Orchestration}
A dedicated dashboard state layer coordinates data fetching and updates. From the received feed data, the dashboard derives indicators such as online feeds, total people, average wait time, and attention counts, then exposes them to the widgets.

\subsection{Live Update Reliability}
The live-update layer includes:
\begin{itemize}
    \item reconnection delays,
    \item a heartbeat check,
    \item a guard against stale data,
    \item fallback polling when live updates fail,
    \item grouped updates to keep the interface responsive.
\end{itemize}
When the connection drops, the dashboard falls back to periodic polling instead of freezing, so operators still have a usable view during transient issues.

\subsection{Live Video Delivery}
Live feed cards use a transport decision layer to choose the most suitable viewing mode. For each running feed, the frontend checks whether live viewing is ready. If the higher-quality live stream is available, it is used first; otherwise the dashboard falls back to the simpler video stream.

The live-streaming flow is handled through the browser's live video connection support \cite{webrtc_w3c}:
\begin{itemize}
    \item create a connection for receiving video,
    \item gather the network information needed for signaling,
    \item send the browser offer to the backend,
    \item apply the backend response and switch the feed card to live state.
\end{itemize}

Reconnection uses increasing delays between attempts. When the problem is not recoverable, the system stops retrying and switches to the fallback stream so the dashboard remains usable.

When server-side annotations are active, the dashboard hides local overlays to avoid drawing the same boxes twice on the same video.

\subsection{Dashboard UI Integration}
The dashboard page brings together live feed controls, setup dialogs, source onboarding, KPI visualizations, and charts. In this sprint, the emphasis was on making those elements work together as one operational interface.

\subsection{Integration Testing}
Frontend integration behavior is validated through scenario-based tests that cover dashboard setup, live event handling, transport fallback, feed routing, and error behavior.

\section{Conclusion}
By the end of Sprint 5, the dashboard could react to live state, switch transport modes, and stay usable when the live connection failed. Sprint 6 then built on that working interface for alert automation.
